\documentclass[a4paper,11pt]{article}
\def\email#1{{\tt#1}}

\def\N{\mbox{\rm{I}\hspace{-.15em}\rm{N}}}
\def\Z{\mbox{\rm{Z}\hspace{-.28em}\rm{Z}}}
\def\RR{\mbox{\rm{I}\hspace{-.15em}\rm{R}}}
\def\C{\mbox{\rm{l}\hspace{-.50em}\rm{C}}}
\def\Q{\mbox{\rm{l}\hspace{-.50em}\rm{Q}}}
\def\F{\mbox{\rm{I}\hspace{-.20em}\rm{F}}}

\usepackage{amssymb,amsmath}
\usepackage{url}
\usepackage{fullpage}
\begin{document}

\newtheorem{theo}{Theorem}
\newtheorem{coro}{Corollary}[theo]
\newtheorem{lemme}{Lemma}
\renewcommand{\thecoro}{\thetheo.\arabic{coro}}



\parindent=0cm


\section*{\center IN100 Initiation au Python\\
Contrôle continu 2}

\vspace{-0.3 cm}

\begin{table}[h!]
\centering
\renewcommand{\arraystretch}{1.4} % augmente la hauteur des lignes
\begin{tabular}{|p{10cm}|p{3cm}|p{2cm}|}
\hline
\textbf{Étudiant :} & \textbf{Heure Début }  & \textbf{Note (/5) }  \\
& & \\
\hline
\textbf{Correcteur :} & \textbf{Heure Fin } & \\
& & \\
\hline

\hline
\end{tabular}
\end{table}
\vspace{-0.4 cm}
\paragraph*{Consignes:} Vous avez 20 minutes pour coder la solution. Seules les types et les instructions vues en cours sont autorisées. Tout usage d'une fonction 
venant d'une bibliothèque qui n'a pas été vue en cours ne sera pris en compte. 

\vspace{0.2 cm}
\paragraph*{Sujet:}

On considère deux suites $u_n$ et $v_n$ définies par :

\[
\begin{cases}
u_0 = 1 \\
v_0 = 2
\end{cases}
\quad \text{et pour tout } n \ge 1,
\quad
\begin{cases}
u_{n} = 2 u_{n-1} + \frac{1}{3}v_{n-1} \\
v_{n} = \frac{1}{2} u_{n-1} + 4 v_{n-1}.
\end{cases}
\]

\begin{enumerate}
 \item Écrire une fonction Python \texttt{suites(N)} qui renvoie une liste contenant tous les tuples  $(u_n,v_n)$, pour tous les $n\leq N$.  Testez votre fonction sur $N=2$ et $N=4$. 
 \item Écrire une fonction Python \texttt{max\_suite\_u(N)} qui renvoie le maximum de la liste 
 $$[u_1-u_0,u_2-u_1, \ldots , u_N-u_{N-1}].$$ 
 Testez cette fonction sur $N=5$ et $N=10$. 
 \end{enumerate}


\vspace{0.5 cm}
\begin{table}[h!]
\centering
\renewcommand{\arraystretch}{1.4} % augmente la hauteur des lignes
\begin{tabular}{|p{15cm}|}
\hline
\textbf{Remarques étudiant :}    \\
 \\
\hline
\textbf{Remarques correcteur :}  \\
 \\
\hline
\end{tabular}
\end{table}


\newpage

\section*{\center IN100 Initiation au Python\\
Contrôle continu 2}

\vspace{-0.3 cm}

\begin{table}[h!]
\centering
\renewcommand{\arraystretch}{1.4} % augmente la hauteur des lignes
\begin{tabular}{|p{10cm}|p{3cm}|p{2cm}|}
\hline
\textbf{Étudiant :} & \textbf{Heure Début }  & \textbf{Note (/5) }  \\
& & \\
\hline
\textbf{Correcteur :} & \textbf{Heure Fin } & \\
& & \\
\hline

\hline
\end{tabular}
\end{table}

\vspace{0.4 cm}
\paragraph*{Consignes:} Vous avez 20 minutes pour coder la solution. Seules les types et les instructions vues en cours sont autorisées. Tout usage d'une fonction 
venant d'une bibliothèque qui n'a pas été vue en cours ne sera pris en compte. 
\vspace{0.4 cm}

\textbf{Sujet :}

Un relevé topographique est représenté par une \textbf{matrice d’entiers}, dans laquelle chaque
entier correspond à une altitude mesurée sur une grille régulière.
On appelle \textbf{pic local} une case de la matrice dont l’altitude est
\textbf{strictement supérieure} à celle de ses quatre voisins immédiats
(haut, bas, gauche et droite).

\medskip

\textit{Remarque :}  
Les éléments situés sur les bords de la matrice ne peuvent pas être des pics locaux,
car ils ne possèdent pas quatre voisins.

\medskip

\textit{Exemple :}
\[
\begin{pmatrix}
10 & 12 & 11 & 9 \\
14 & 18 & 13 & 8 \\
9  & 15 & 10 & 7 \\
6  & 8  & 5  & 4
\end{pmatrix}.
\]

Dans cette matrice, la valeur \(18\) (ligne 1, colonne 1) est un pic local car
\(18 > 14\), \(18 > 12\), \(18 > 13\) et \(18 > 15\).

\medskip

\textbf{Questions :}

\begin{enumerate}
\item Écrire une fonction \texttt{est\_pic(matrice, i, j)} qui prend en argument une matrice
d’entiers et deux indices \(i\) et \(j\), et qui renvoie \texttt{True} si l’élément situé à la
ligne \(i\) et à la colonne \(j\) est un pic local, et \texttt{False} sinon.

\item Écrire une fonction \texttt{compter\_pics(matrice)} qui prend en argument une matrice
d’entiers représentant le relevé topographique et qui renvoie le nombre total de pics locaux
présents dans cette matrice.
\end{enumerate}




%\vspace{0.5 cm}
\begin{table}[h!]
\centering
\renewcommand{\arraystretch}{1.4} % augmente la hauteur des lignes
\begin{tabular}{|p{15cm}|}
\hline
\textbf{Remarques étudiant :}    \\
 \\
\hline
\textbf{Remarques correcteur :}  \\
 \\
\hline
\end{tabular}
\end{table}

\newpage

\section*{\center IN100 Initiation au Python\\
Contrôle continu 2}

\vspace{0.3 cm}

\begin{table}[h!]
\centering
\renewcommand{\arraystretch}{1.4} % augmente la hauteur des lignes
\begin{tabular}{|p{10cm}|p{3cm}|p{2cm}|}
\hline
\textbf{Étudiant :} & \textbf{Heure Début }  & \textbf{Note (/5) }  \\
& & \\
\hline
\textbf{Correcteur :} & \textbf{Heure Fin } & \\
& & \\
\hline

\hline
\end{tabular}
\end{table}
\vspace{-0.4 cm}
\paragraph*{Consignes:} Vous avez 20 minutes pour coder la solution. Seules les types et les instructions vues en cours sont autorisées. Tout usage d'une fonction 
venant d'une bibliothèque qui n'a pas été vue en cours ne sera pris en compte. 



\paragraph*{Sujet:}


Nous allons écrire un programme permettant de faire le suivi des pas journaliers. 
On dispose d'une liste \texttt{pas} d'entiers, o\`u \texttt{pas[i]} est le nombre de pas effectu\'es le jour \texttt{i}.\\
Exemple :
\begin{verbatim}
pas = [3200, 5100, 4900, 8000, 8000, 2000, 6500]
\end{verbatim}

\begin{enumerate}
\item \'Ecrire une fonction \texttt{jours\_actifs(pas, seuil)} qui renvoie la liste des indices des jours o\`u
\texttt{pas[i] >= seuil}.\\

\vspace{0.1 cm}

\textit{Exemple :} \texttt{jours\_actifs(pas, 6000)} renvoie \texttt{[3, 4, 6]}.

\item \'Ecrire une fonction \texttt{plus\_longue\_serie(pas, seuil)} qui renvoie la longueur de la plus longue suite d'indices
\textbf{cons\'ecutifs} dans la liste calculée par la fonction de la question 1. 

\vspace{0.1 cm}
\textit{Exemple :} ici \texttt{J = [3,4,6]}, donc la plus longue s\'erie vaut 2 (3 puis 4).
\end{enumerate}


\vspace{0.5 cm}
\begin{table}[h!]
\centering
\renewcommand{\arraystretch}{1.4} % augmente la hauteur des lignes
\begin{tabular}{|p{15cm}|}
\hline
\textbf{Remarques étudiant :}    \\
 \\
\hline
\textbf{Remarques correcteur :}  \\
 \\
\hline
\end{tabular}
\end{table}

\newpage


\section*{\center IN100 Initiation au Python\\
Contrôle continu 2}


\begin{table}[h!]
\centering
\renewcommand{\arraystretch}{1.4} % augmente la hauteur des lignes
\begin{tabular}{|p{10cm}|p{3cm}|p{2cm}|}
\hline
\textbf{Étudiant :} & \textbf{Heure Début }  & \textbf{Note (/5) }  \\
& & \\
\hline
\textbf{Correcteur :} & \textbf{Heure Fin } & \\
& & \\
\hline

\hline
\end{tabular}
\end{table}

\vspace{-0.3 cm}

\paragraph*{Consignes:} Vous avez 20 minutes pour coder la solution. Seules les types et les instructions vues en cours sont autorisées. Tout usage d'une fonction 
venant d'une bibliothèque qui n'a pas été vue en cours ne sera pris en compte. 

\paragraph*{Sujet:} 
Durant un cours de sport, un professeur souhaite organiser un tournoi de tennis de table. Sa classe comporte 26 élèves, qu'il numérote de 1 à 26. Il dispose de 13 tables sur lesquels les élèves s'affrontent. On représente les terrains par une liste de tables, et chaque table par un tuple constituant une paire d'élèves. Par exemple, une disposition de départ des élèves est la suivante: 
\begin{center}
\texttt{[(1,2),(3,4),(5,6),...,(23,24),(25,26)]}
\end{center}


\begin{enumerate}
  \item[1.] Une fois un match terminé sur une table, on souhaite que le premier élèves du tuples correspondant soit le vainqueur. Écrire une fonction Python \texttt{match} qui prend en entrée un tuple d'élèves, et choisit un gagnant au hasard et renvoie un tuple adapté au résultat du match.
  \item[2.] Lorsqu'un élève gagne, il monte à la table voisine d'indice supérieur. Lorsqu'un élève perd, il descend à la table voisine d'indice inférieur. Le perdant de la table d'indice $0$ ne change pas de table, et il en va de même pour le gagnant de la table d'indice $12$.\\   
  Écrire une fonction Python \texttt{round} qui prend en entrée une liste de tuples représentant les tables, et qui simule un match sur chaque table et renvoie la liste des tables après que les élèves se soient déplacés.
\end{enumerate}
 
Pour générer aléatoirement des nombres entiers, nous rappelons qu'on peut utiliser la fonction \texttt{random.randint(a,b)} qui renvoie un nombre entier entre $a$ et $b$. 

\vspace{0.2 cm}
Exemple d'usage:
%\vspace{-0.2 cm}
\begin{verbatim}
	import random	
	print(random.randint(2,15))
\end{verbatim}


\vspace{-0.6 cm}
\begin{table}[h!]
\centering
\renewcommand{\arraystretch}{1.4} % augmente la hauteur des lignes
\begin{tabular}{|p{15cm}|}
\hline
\textbf{Remarques étudiant :}    \\
 \\
\hline
\textbf{Remarques correcteur :}  \\
 \\
\hline
\end{tabular}
\end{table}


\newpage


\section*{\center IN100 Initiation au Python\\
Contrôle continu 2}

\vspace{0.5 cm}

\begin{table}[h!]
\centering
\renewcommand{\arraystretch}{1.4} % augmente la hauteur des lignes
\begin{tabular}{|p{10cm}|p{3cm}|p{2cm}|}
\hline
\textbf{Étudiant :} & \textbf{Heure Début }  & \textbf{Note (/5) }  \\
& & \\
\hline
\textbf{Correcteur :} & \textbf{Heure Fin } & \\
& & \\
\hline

\hline
\end{tabular}
\end{table}

\paragraph*{Consignes:} Vous avez 20 minutes pour coder la solution. Seules les types et les instructions vues en cours sont autorisées. Tout usage d'une fonction 
venant d'une bibliothèque qui n'a pas été vue en cours ne sera pris en compte. 



\paragraph*{Sujet:}

On considère un flux de messages représenté par une liste de chaînes de caractères.  
Chaque message est associé à un niveau de priorité : un message \emph{urgent} commence par le caractère \texttt{!},
 tandis qu'un message \emph{non urgent} ne commence pas par \texttt{!}.


Le but est de séparer ces messages en deux \textbf{files}, une pour les messages urgents et l'autre pour les messages non-urgents. 
Pour cela, on utilise :
\begin{itemize}
    \item une liste \texttt{messages} représentant une file de message, dans l'ordre de leur arrivée. 
   \item deux \textbf{files}, une pour les messages urgents et l'autre pour les messages non-urgents. 
\end{itemize}

\begin{enumerate}
    \item Écrire une fonction Python \texttt{enfiler(f, x)} qui ajoute le message  \texttt{x} dans la file  \texttt{f} et une fonction \texttt{defiler(f)} qui enlève un élement de la file \texttt{f}. 
    \item Écrire une fonction Python \texttt{traiter\_messages(messages)} qui traite les messages de la file \texttt{messages} un par un et :
    \begin{itemize}
        \item ajoute le message (sans le caractère \texttt{!}) dans la file \texttt{urgents} s'il est urgent,
        \item ajoute le message dans la file \texttt{non\_urgents} sinon.
    \end{itemize}
Vérifier que tous les messages ont été traités et que la file \texttt{messages} est vide à la fin du processus.
\end{enumerate}



\begin{table}[h!]
\centering
\renewcommand{\arraystretch}{1.4}
\begin{tabular}{|p{15cm}|}
\hline
\textbf{Remarques étudiant :} \\
 \\
\hline
\textbf{Remarques correcteur :} \\
 \\
\hline
\end{tabular}
\end{table}


\end{document}
